\documentclass[a4paper]{article}

\usepackage[utf8]{inputenc}
\usepackage{amsmath}
\usepackage{graphicx}
\usepackage[colorinlistoftodos]{todonotes}

\usepackage{multicol}
\usepackage{makeidx}
\usepackage{hyperref}
\usepackage{caption}
\usepackage{amsfonts}
\usepackage{amssymb}
\usepackage[utf8]{inputenc}
\usepackage{verbatim}
\usepackage{listings}
\lstset{language=C++, showstringspaces=false, tabsize=2, breaklines=true, title=\lstname}

\usepackage[margin=0.5in]{geometry}

\title{TP2 Metnum}

\author{Pedro Rodriguez, Gonzalo, Ezequiel}

\date{\today}

\makeindex

\begin{document}
\newgeometry{margin=2cm}
\pagenumbering{gobble}
\raggedleft
\includegraphics[width=8cm]{logo1.jpg}\\

\raggedright
\vspace{3cm}
{\Huge \bfseries Trabajo Práctico 1}
\rule{\textwidth}{0.02in}
\large Martes 16 de Septiembre de 2014 \hfill Sistemas Operativos
\vspace{1.5cm}

 
\centering \LARGE Scheduling 
\vspace{1.5cm}

\normalsize
\begin{tabular}{|l@{\hspace{4ex}}c@{\hspace{4ex}}l|}
        \hline
        \rule{0pt}{1.2em}Integrante & LU & Correo electr\'onico\\[0.2em]
        \hline
        \rule{0pt}{1.2em}  Pedro Rodriguez &197/12 &\tt pedro3110.jim@gmail.com \\[0.2em]
        \rule{0pt}{1.2em}  Gonzalo Benegas &??? &\tt  gsbenegas@gmail.com \\[0.2em]
        \rule{0pt}{1.2em}  Ezequiel  Barrios &??? &\tt ezequiel.barrios@gmail.com \\[0.2em]
        \hline
\end{tabular}

\vspace{1.0cm}
\raggedright

\begin{multicols}{2}
\includegraphics[width=8cm]{logo-uba.png}

\columnbreak
\vspace*{4.5cm}
\raggedleft
\textbf{Facultad de Ciencias Exactas y Naturales}\\
Universidad de Buenos Aires\\
\small
Ciudad Universitaria - (Pabellon I/Planta Baja)\\
Intendente G\"uiraldes 2160 - C1428EGA\\
Ciudad Autonoma de Buenos Aires - Rep. Argentina\\
Tel/Fax: (54 11) 4576-3359\\
http://www.fcen.uba.ar
\end{multicols}

\restoregeometry

\clearpage

\pagenumbering{arabic}

\tableofcontents

\vspace{3cm}

\begin{abstract}
En el presente trabajo pr\'actico estudiaremos algunos de los m\'etodos m\'as com\'unmente utilizados 
en la actualidad por distintos sistemas operativos para manejar correcta y eficientemente los 
diversos procesos que se ejecutan concurrentemente en m\'aquinas con uno o m\'as procesadores. \\
Intentaremos detectar las ventajas y desventajas de cada m\'etodo, as\'i como los escenarios en los
cuales uno puede ser m\'as eficiente que otro. Para esto, dividimos el TP en tres partes: en la
primera, presentamos el simulador simusched y lo corremos para algunas tareas espec\'ificas. En la
segunda parte, extendemos el simulador con nuevos schedulers implementados por nosotros y finalmente,
en la tercera parte evaluamos y analizamos los distintos algoritmos de scheduling ya presentados. \\
\end{abstract}

\index{Parte I - Entendiendo el simulador simusched}
\section{Parte I - Entendiendo el simulador simusched}

\subsection{Ejercicio 1}
En este ejercicio, programamos la tarea TaskConsola, que simular\'ia una tarea interactiva con
el usuario. \\
Para la implementaci\'on de esta tarea, primero tuvimos que registrarla para que el simulador la 
reconozca en el archivo tasks.cpp, con la funci\'on tasksinit(void). En ese mismo archivo implementamos
la tarea, que para el proceso n\'umero pid recibe los tres par\'ametros: n, bmin y bmax. Simplemente
hacemos un ciclo que cicle n veces y que cada vez tome un entero al azar $r$ entre $bmin$ y 
$bmax$, y cada vez llamamos a la funci\'on usoIO(pid, r), que simula hacer uso de dispositivos de
entrada/salida. \\
Nota(1): se denomina llamada bloqueante a aquellas llamadas en las que, si lo que esta  
solicita no est\'a disponible, entonces el proceso llamador se queda bloqueado a la espera de un
resultado. Es decir, no permite que otros procesos dependientes de este sigan ejecutando. Por 
el contrario, en una llamada no bloqueante, si el proceso llamador no encuentra lo que estaba
buscando, el proceso devuelve de todas formas alg\'un resultado, permitiendo que otros procesos
sigan ejecutando sin tener que esperar a que este proceso llamador reciba el resultado que est\'a
esperando. \\
Nota(2): para elegir el n\'umero pseudo-aleatorio hacemos uso de la funci\'on rand() provista
por la librer\'a standard de C. \\

\subsection{Ejercicio 2}
Para seguir entendiendo el funcionamiento del simulador simusched, ahora pasamos a escribir un lote
de 3 tareas distintas: una intensiva en CPU y las otras dos de tipo interactivo. Vamos a ejecutar y
graficar la simulaci\'on usando el algoritmo FCFS para 1, 2 y 3 n\'ucleos. \\


\index{Parte II - Extendiendo el simulador con nuevos schedulers}
\section{Parte II - Extendiendo el simulador con nuevos schedulers}

\subsection{Ejercicio 3}
\subsection{Ejercicio 4}
\subsection{Ejercicio 5}
En este ejercicio, basados en el paper de Waldspurger, C.A. y Weihl en el cu\'al presentan un 
novedoso sistema de Scheduling llamado $Lottery Scheduling$, implementamos una clase que llamamos
SchedLottery que recibe como par\'ametros el quantum que se va a estar simulando y una semilla 
pseudo-aleatoria, que el scheduler va a utilizar para el manejo de los procesos. Como dice en el
enunciado, b\'asicamente nos interes\'o implementar la idea b\'asica del algoritmo y la 
optimizaci\'on de tickets compensatorios. \\




\index{Parte III - Evaluando los algoritmos de scheduling}
\section{Parte III - Evaluando los algoritmos de scheduling}

\subsection{Ejercicio 6}
\subsection{Ejercicio 7}
\subsection{Ejercicio 8}
\subsection{Ejercicio 9}
\subsection{Ejercicio 10}














\end{document}
